\documentclass[10pt,a4paper]{amsart}
\usepackage{amsmath,amssymb,amsthm,graphicx}
% Retained preamble for the Stacks Project chapter divisors.tex (Divisors).
% The chapter file pinned-original-source.tex carries only \input{preamble} on its first line, so no
% preamble text lives inside the chapter itself.  The Stacks Project uses one single shared preamble for
% every chapter; the snapshot stages it verbatim in sources/stacks-chow/preamble.tex.  The block below is
% copied mechanically, line for line, from that file's lines 14-18 (the Xy-pic setup) and lines 154-217
% (the native theorem declarations together with the mathematical macros).  Nothing was added, removed or
% redefined.  The wrapper class and the amsmath/amssymb/amsthm/graphicx base packages are supplied by the
% assembler, so the class/\externaldocument part of the shared preamble is not repeated here.
% For commutative diagrams we use Xy-pic
\usepackage[all]{xy}

% We use 2cell for 2-commutative diagrams.
\xyoption{2cell}
% Theorem environments.
%
\theoremstyle{plain}
\newtheorem{theorem}[subsection]{Theorem}
\newtheorem{proposition}[subsection]{Proposition}
\newtheorem{lemma}[subsection]{Lemma}

\theoremstyle{definition}
\newtheorem{definition}[subsection]{Definition}
\newtheorem{example}[subsection]{Example}
\newtheorem{exercise}[subsection]{Exercise}
\newtheorem{situation}[subsection]{Situation}

\theoremstyle{remark}
\newtheorem{remark}[subsection]{Remark}
\newtheorem{remarks}[subsection]{Remarks}

\numberwithin{equation}{subsection}

% Macros
%
\def\lim{\mathop{\mathrm{lim}}\nolimits}
\def\colim{\mathop{\mathrm{colim}}\nolimits}
\def\Spec{\mathop{\mathrm{Spec}}}
\def\Hom{\mathop{\mathrm{Hom}}\nolimits}
\def\Ext{\mathop{\mathrm{Ext}}\nolimits}
\def\SheafHom{\mathop{\mathcal{H}\!\mathit{om}}\nolimits}
\def\SheafExt{\mathop{\mathcal{E}\!\mathit{xt}}\nolimits}
\def\Sch{\mathit{Sch}}
\def\Mor{\mathop{\mathrm{Mor}}\nolimits}
\def\Ob{\mathop{\mathrm{Ob}}\nolimits}
\def\Sh{\mathop{\mathit{Sh}}\nolimits}
\def\NL{\mathop{N\!L}\nolimits}
\def\CH{\mathop{\mathrm{CH}}\nolimits}
\def\proetale{{pro\text{-}\acute{e}tale}}
\def\etale{{\acute{e}tale}}
\def\QCoh{\mathit{QCoh}}
\def\Ker{\mathop{\mathrm{Ker}}}
\def\Im{\mathop{\mathrm{Im}}}
\def\Coker{\mathop{\mathrm{Coker}}}
\def\Coim{\mathop{\mathrm{Coim}}}

% Boxtimes
%
\DeclareMathSymbol{\boxtimes}{\mathbin}{AMSa}{"02}

%
% Macros for moduli stacks/spaces
%
\def\QCohstack{\mathcal{QC}\!\mathit{oh}}
\def\Cohstack{\mathcal{C}\!\mathit{oh}}
\def\Spacesstack{\mathcal{S}\!\mathit{paces}}
\def\Quotfunctor{\mathrm{Quot}}
\def\Hilbfunctor{\mathrm{Hilb}}
\def\Curvesstack{\mathcal{C}\!\mathit{urves}}
\def\Polarizedstack{\mathcal{P}\!\mathit{olarized}}
\def\Complexesstack{\mathcal{C}\!\mathit{omplexes}}
% \Pic is the operator that assigns to X its picard group, usage \Pic(X)
% \Picardstack_{X/B} denotes the Picard stack of X over B
% \Picardfunctor_{X/B} denotes the Picard functor of X over B
\def\Pic{\mathop{\mathrm{Pic}}\nolimits}
\def\Picardstack{\mathcal{P}\!\mathit{ic}}
\def\Picardfunctor{\mathrm{Pic}}
\def\Deformationcategory{\mathcal{D}\!\mathit{ef}}
\usepackage{cleveref}
\crefname{originallemma}{original lemma}{original lemmas}
\Crefname{originallemma}{Original lemma}{Original lemmas}
\makeatletter
\newlabel{lemma-integral-effective-Cartier-divisor-dvr}{{lemma-integral-effective-Cartier-divisor-dvr}{source}{}{sourcepointer-1}{}}
\newlabel{lemma-integral-effective-Cartier-divisor-dvr@cref}{{[originallemma][0][]lemma-integral-effective-Cartier-divisor-dvr}{[1][0][]source}{}{}{}}
\makeatother
\crefname{originallemma}{original lemma}{original lemmas}
\Crefname{originallemma}{Original lemma}{Original lemmas}
\makeatletter
\newlabel{lemma-weil-divisor-is-cartier-UFD}{{lemma-weil-divisor-is-cartier-UFD}{source}{}{sourcepointer-2}{}}
\newlabel{lemma-weil-divisor-is-cartier-UFD@cref}{{[originallemma][0][]lemma-weil-divisor-is-cartier-UFD}{[1][0][]source}{}{}{}}
\makeatother
\crefname{originalreference}{original reference}{original references}
\Crefname{originalreference}{Original reference}{Original references}
\makeatletter
\newlabel{more-algebra-lemma-regular-local-UFD}{{more-algebra-lemma-regular-local-UFD}{source}{}{sourcepointer-3}{}}
\newlabel{more-algebra-lemma-regular-local-UFD@cref}{{[originalreference][0][]more-algebra-lemma-regular-local-UFD}{[1][0][]source}{}{}{}}
\makeatother
\crefname{originalreference}{original reference}{original references}
\Crefname{originalreference}{Original reference}{Original references}
\makeatletter
\newlabel{topology-definition-codimension}{{topology-definition-codimension}{source}{}{sourcepointer-4}{}}
\newlabel{topology-definition-codimension@cref}{{[originalreference][0][]topology-definition-codimension}{[1][0][]source}{}{}{}}
\makeatother
\title{A closed subscheme covered by effective Cartier divisors and a codimension two closed subset is, away from codimension two, an effective Cartier divisor supported on it}
\author{The Stacks Project Authors}
\date{stacks-project:divisors.tex:local-snapshot-20261012; GNU FDL 1.2 or later}
\begin{document}
\maketitle
\noindent\textit{Editorial note.} All delivered mathematical text is the author's own, copied from the pinned Stacks Project chapter file pinned-original-source.tex; no statement or proof step was rewritten, reordered or removed.\par\medskip
\noindent\textit{Editorial note.} Following the multi-claim convention for this source, source\_url is the source-level URL https://stacks.math.columbia.edu/ and the source label lemma-codim-1-part is carried in claim\_key and here rather than in a per-claim tag URL.\par\medskip
\noindent\textit{Editorial note.} The preamble is the chapter's shared Stacks Project preamble, retained as sources/stacks-divisors/retained-preamble.tex: a line-for-line excerpt of the shared Stacks preamble staged with this snapshot. No macro was invented or redefined, and this card adds no preamble text of its own.\par\medskip
\noindent\textit{Editorial note.} Statement span 3075-3100 and proof span 3102-3155 come from the halfA pool card operations/pool-cards/stacks-divisors-halfA.json and were re-read in the pinned chapter file before this card was written.\par\medskip
\noindent\textit{Editorial note.} The proof invokes two equivalents of the same fact about discrete valuation rings and uses the phrase 'Proof of the claim.' as an internal heading; both are the author's text and were kept verbatim.\par\medskip
\noindent\textit{Editorial note.} 4 cross-references occur inside the counted spans; they are listed in external\_references as navigation pointers only.\par\medskip
\noindent\textit{Allowed standard prerequisites (editorial).}\begin{itemize}
\item An integral effective Cartier divisor has a discrete valuation ring at its generic point, and one can speak of the valuation of an element of the stalk of an ideal sheaf there. Cited in the same chapter at proof lines 3104-3106. Supplies the local rings whose valuations define the multiplicities a\_i.
\item Local equation of an effective Cartier divisor at a point of the divisor, so that the divisors D\_i through a point are cut out by prime elements of the local ring. Used at proof lines 3114-3116. Produces the prime elements f\_i in terms of which the product of the f\_i to the a\_i is formed.
\item A regular local ring is a unique factorisation domain, cited from the chapter on more algebra, and the same-chapter statement that irreducible codimension one components of a closed subscheme in a UFD situation are effective Cartier divisors. Cited at proof lines 3147-3152. Supplies the existence part of the final sentence: a regular or locally UFD ambient produces the required divisors D\_i.
\end{itemize}\medskip
\ifdefined\hypertarget\hypertarget{sourcepointer-1}{}\fi
\noindent\textit{Source reference.} Original lemma lemma-integral-effective-Cartier-divisor-dvr: sources/stacks-divisors/pinned-original-source.tex, source lines 3006--3006. Same-chapter reference used in the retained author text at source line(s) 3106. It is supplied only as a navigation pointer so that the retained text compiles; the counted item delivers the author's own statement and the author's complete proof, and no part of the referenced item is copied in or substituted for the counted argument.\par\medskip
\ifdefined\hypertarget\hypertarget{sourcepointer-2}{}\fi
\noindent\textit{Source reference.} Original lemma lemma-weil-divisor-is-cartier-UFD: sources/stacks-divisors/pinned-original-source.tex, source lines 3056--3056. Same-chapter reference used in the retained author text at source line(s) 3153, 3520, 6926. It is supplied only as a navigation pointer so that the retained text compiles; the counted item delivers the author's own statement and the author's complete proof, and no part of the referenced item is copied in or substituted for the counted argument.\par\medskip
\ifdefined\hypertarget\hypertarget{sourcepointer-3}{}\fi
\noindent\textit{Source reference.} Original reference more-algebra-lemma-regular-local-UFD: sources/stacks-divisors/pinned-original-source.tex, source lines 3149--3149. Same-chapter reference used in the retained author text at source line(s) 3149, 3570. It is supplied only as a navigation pointer so that the retained text compiles; the counted item delivers the author's own statement and the author's complete proof, and no part of the referenced item is copied in or substituted for the counted argument.\par\medskip
\ifdefined\hypertarget\hypertarget{sourcepointer-4}{}\fi
\noindent\textit{Source reference.} Original reference topology-definition-codimension: sources/stacks-divisors/pinned-original-source.tex, source lines 3084--3084. Same-chapter reference used in the retained author text at source line(s) 3084. It is supplied only as a navigation pointer so that the retained text compiles; the counted item delivers the author's own statement and the author's complete proof, and no part of the referenced item is copied in or substituted for the counted argument.\par\medskip
\begin{lemma}
\label{lemma-codim-1-part}
Let $X$ be a Noetherian scheme. Let $Z \subset X$ be a closed subscheme.
Assume there exist
\begin{enumerate}
\item a collection of integral effective Cartier divisors
$D_i \subset X$, $i \in I$
\item a closed subset $Z' \subset X$ all of whose irreducible components
have codimension $\geq 2$ in $X$
(Topology, Definition \ref{topology-definition-codimension})
\end{enumerate}
such that $Z \subset Z' \cup \bigcup_{i \in I} D_i$ set-theoretically.
Then there exist integers $a_i \geq 0$ with $a_i = 0$ for almost all $i \in I$
such that the effective Cartier divisor
$$
D = \sum a_i D_i
$$
is contained in $Z$ and such that the inclusion morphism $D \to Z$
is an isomorphism away from codimension $2$ in $X$
(in the sense that there exists an open $U \subset Z$ such that
$D \cap U \to Z \cap U$ is an isomorphism and such that every
irreducible component of $Z \setminus U$ has codimension $\geq 2$ in $X$).
When $Z$ is nowhere dense in $X$ existence of the $D_i$, $i \in I$ and $Z'$
is guaranteed if $\mathcal{O}_{X, x}$ is a UFD for all $x \in Z$
or if $X$ is regular.
\end{lemma}

\begin{proof}
Let $\xi_i \in D_i$ be the generic point and let
$\mathcal{O}_i = \mathcal{O}_{X, \xi_i}$ be the local ring
which is a discrete valuation ring by
Lemma \ref{lemma-integral-effective-Cartier-divisor-dvr}.
Let $a_i \geq 0$ be the minimal valuation of an element of
$\mathcal{I}_{Z, \xi_i} \subset \mathcal{O}_i$.
Of course $a_i > 0$ only if $D_i$ is an irreducible component
of $Z$ and hence $a_i > 0$ only for a finite number of $i \in I$.
We claim that the effective Cartier divisor $D = \sum a_i D_i$ works.

\medskip\noindent
Namely, suppose that $x \in X$. Let $A = \mathcal{O}_{X, x}$.
Let $D_1, \ldots, D_n$ be the pairwise distinct divisors
$D_i$ such that $x \in D_i$ and $a_i > 0$.
For $1 \leq i \leq n$ let $f_i \in A$ be a local equation for $D_i$.
Then $f_i$ is a prime element of $A$ and $\mathcal{O}_i = A_{(f_i)}$. Let
$I = \mathcal{I}_{Z, x} \subset A$ be the stalk of the
ideal sheaf of $Z$. By our choice of $a_i$ we have
$I A_{(f_i)} = f_i^{a_i}A_{(f_i)}$. We claim that
$I \subset (\prod_{i = 1, \ldots, n} f_i^{a_i})$.

\medskip\noindent
Proof of the claim. The localization map
$\varphi : A/(f_i) \to A_{(f_i)}/f_iA_{(f_i)}$ is injective as
the prime ideal $(f_i)$ is the inverse image of the maximal ideal
$f_iA_{(f_i)}$. By induction on $n$ we deduce that
$\varphi_n : A/(f_i^n)\to A_{(f_i)}/f_i^nA_{(f_i)}$ is also injective.
Since $\varphi_{a_i}(I) = 0$, we have $I \subset (f_i^{a_i})$.
Thus, for any $x \in I$, we may write $x = f_1^{a_1}x_1$
for some $x_1 \in A$. Since $D_1, \ldots, D_n$ are pairwise
distinct, $f_i$ is a unit in $A_{(f_j)}$ for $i \not = j$.
Comparing $x$ and $x_1$ at $A_{(f_i)}$
for $n \geq i > 1$, we still have $x_1 \in (f_i^{a_i})$.
Repeating the previous process, we inductively write
$x_i = f_{i + 1}^{a_{i + 1}}x_{i + 1}$ for any $n > i \geq 1$.
In conclusion, $x \in (\prod_{i = 1, \ldots n} f_i^{a_i})$
for any $x \in I$ as desired.

\medskip\noindent
The claim shows that $\mathcal{I}_Z \subset \mathcal{I}_D$, i.e., that
$D \subset Z$. Moreover, we also see that $D$ and $Z$ agree at the $\xi_i$,
which proves that $D \to Z$ is an isomorphism away from codimension $2$ on $X$.

\medskip\noindent
To see the final statements we argue as follows. A regular local
ring is a UFD (More on Algebra, Lemma
\ref{more-algebra-lemma-regular-local-UFD}) hence it suffices
to argue in the UFD case. In that case, let
$D_i$ be the irreducible components of $Z$
which have codimension $1$ in $X$.
By Lemma \ref{lemma-weil-divisor-is-cartier-UFD} each $D_i$
is an effective Cartier divisor.
\end{proof}

\end{document}
